\subsection{Блоки}
%%%%%%%%%%%%%%%%%%%%%%%%%%%%%%%%%%%%%%%%%%%%%%%%%%%%

 \z {Вычислите ускорения грузиков массой $m_2$:
 		\begin{center}%
 			\includegraphics[width=0.65\textwidth]{4.6.1}%
		\end{center}%
 }
%
{\puan $ a_2=\frac{4m_1+m_2}{16m_1+m_2}g;$~\puan $ a_2=\frac{m_2-2m_1}{4m_1+m_2}g;$~\puan $ a_2=\frac{m_2-4m_1}{16m_1+m_2}g$}
%%%%%%%%%%%%%%%%%%%%%%%%%%%%%%%%%%%%%%%%%%%%%%%%%%%% 
 \z {Вычислите ускорения грузиков массой $m_2$:
 		\begin{center}%
 			\includegraphics[width=0.6\textwidth]{4.6.2}%
		\end{center}%
 }
%
{\puan $ a_2=\frac{m_2-m_1}{m_1+m_2}g;$~\puan $ a_2=a_1=g$}
%%%%%%%%%%%%%%%%%%%%%%%%%%%%%%%%%%%%%%%%%%%%%%%%%%%% 

 \z {Вычислите ускорения грузиков массой $m_3$:
 		\begin{center}%
 			\includegraphics[width=0.8\textwidth]{4.6.5}%
		\end{center}%
 }
%
{\puan $ a_3=3\frac{4m_1-2 m_2+3 m_3}{16m_1+4m_2+9m_3}g;$~\puan $ a_3=3\frac{4m_1-4 m_2+3 m_3}{16m_1+4m_2+9m_3}g;$~\puan $ a_3=3\frac{8m_1-2m_2+3 m_3}{64m_1+16m_2+9m_3}g$}
%%%%%%%%%%%%%%%%%%%%%%%%%%%%%%%%%%%%%%%%%%%%%%%%%%%% 
 \z[{\begin{wrapfigure}[5]{r}{0.2\textwidth}%
		\vspace{-15pt}%
		\begin{flushright}%
 			\includegraphics[width=0.15\textwidth]{4.6.3}%
		\end{flushright}%
	\end{wrapfigure}%
	}]%
	{Массы трех грузов в системе равны $M$, правый нижний груз имеет массу $2M$. Верхние грузы вначале удерживают, и в тот момент, когда скорости остальных двух грузов равны нулю, грузы отпускают. Найти ускорения каждого из грузов и натяжение длинной нити.}
%
{$a_\text{верх} = 1.25g;~ a_\text{ниж.лег.}=0.5g;~a_\text{ниж.тяж.}=0.75g$}
%%%%%%%%%%%%%%%%%%%%%%%%%%%%%%%%%%%%%%%%%%%%%%%%%%%%
 \z[{\begin{wrapfigure}[5]{r}{0.2\textwidth}%
		\vspace{-15pt}%
		\begin{flushright}%
 			\includegraphics[width=0.15\textwidth]{4.6.4}%
		\end{flushright}%
	\end{wrapfigure}%
	}]%
	{В системе на рисунке справа маленькие грузы имеют массы m, большой груз массу $3m$. С какой силой груз $m$ давит на $3m$ во время движения?}
%
{$N=0.75mg$}
%%%%%%%%%%%%%%%%%%%%

%%%%%%%%%%%%%%%%%%%%%%%%%%%%%%%%%%%%%%%%%%%%%%%%%%%%